\section{Calibrated curvature and linear coercivity}\label{sec:linear-coercivity}

The K\"ahler value of the curvature pencil admits a stronger estimate than
non-saturation away from that value.  Its curvature model is not small, but
its action on the calibrated spinor vanishes.  Retaining this cancellation
in the covariant Hessian identity removes the background-volume term from
the second-order estimate.  A weak Bianchi test then bounds the spinor
energy linearly by the curvature deficit.  Conformal normalization turns
this estimate into a scale-invariant inequality involving only traceless
Ricci energy and the scalar--Weyl deficit.

\subsection{Actual input and the retained weighted identities}

Let $(M,g)$ be a closed, connected, oriented smooth four-manifold.  Fix a
Spin$^c$ structure whose unperturbed Seiberg--Witten equations are solvable
for every metric conformal to $g$, and let $a\in H^2(M;\mathbb R)$ be its
real first Chern class.  Throughout the argument this one structure is
fixed while the auxiliary conformal weights vary.  In an application
quantified over all metrics, the selected structure may depend on the
original metric.  Set
\begin{equation}\label{eq:lc-energies}
\begin{gathered}
 V=\Vol_g(M),\qquad p=(a_g^+)^2>0,\qquad
 T_\pm=2\chi(M)\pm3\tau(M),\qquad T=T_+,\\
 S_2=\int_M s^2\,d\mu,\qquad
 P=\int_M|\Ric^\circ|^2\,d\mu,\qquad
 R=\int_M|\Ric|^2\,d\mu=P+\frac14S_2,\\
 D=D_2=\frac32 S_2+18\int_M|W^+|^2\,d\mu-72\pi^2p,\qquad
 \Sigma=S_2-32\pi^2p,\qquad Y=(T-a^2)_+.
\end{gathered}
\end{equation}
All unqualified integrals and norms in this section use $g$.  Ricci
curvature has the tensor norm, two-forms have the exterior-algebra norm,
and $W^+$ has the Hilbert--Schmidt norm as an endomorphism of $\Lambda^+$.
The classical scalar estimate and the actual weighted pencil of
Sections~\ref{sec:intrinsic-calibration} and~\ref{sec:weighted-pencil} give
\begin{equation}\label{eq:lc-nonnegative}
 \int_M s_-^2\,d\mu\ge32\pi^2p,\qquad \Sigma\ge0,\qquad D\ge0.
\end{equation}
In particular $R>0$.  No diameter, injectivity-radius, Sobolev,
noncollapse, or full-curvature bound is imposed.

First fix the negative target $\Ric=-g$ and write
\begin{equation}\label{eq:lc-normalized-error}
 r=\frac1V\int_M|\Ric+g|^2\,d\mu,\qquad
 \Xi=rV+D,\qquad \kappa=\frac{\Xi}{V},\qquad
 \Theta=\Xi+\varepsilon V.
\end{equation}
Here $\varepsilon>0$ is a freely selectable allowance in the smooth
weight approximation.  It is chosen after the metric and structure
have been fixed.  Unless explicitly stated otherwise, the weighted
spinor estimates below are used with $\kappa\le1$ and $\varepsilon\le1$.

For clarity we specify the actual solution and the information retained
from the pencil.  Put
\begin{equation}\label{eq:lc-weight}
 w=\sqrt6|W^+|,\qquad q=-s+w,\qquad Q=\max\{q,3\},\qquad
 f_0=\max\{\sqrt Q,s_-/4\}.
\end{equation}
Smooth positive majorants $f$ of $f_0$ can be chosen with arbitrarily
small approximation allowance.  Conformal covariance supplies an actual
solution
\begin{equation}\label{eq:lc-sw}
 D_A\Phi=0,\qquad F_A^+=if\sigma(\Phi),\qquad
 4\nabla_A^*\nabla_A\Phi+(s+f|\Phi|^2)\Phi=0.
\end{equation}
The normalization of the quadratic map is
$|\sigma(\Phi)|^2=|\Phi|^4/8$.  Define
\begin{equation}\label{eq:lc-spinor-notation}
\begin{gathered}
 \rho=|\Phi|^2,\qquad \psi=2\sqrt2\sigma(\Phi),\qquad
 \zeta=\langle\Phi,\nabla_A\Phi\rangle,\qquad U=s+f\rho,\\
 K=\int_M|\nabla_A\Phi|^2\,d\mu,\qquad
 B=\int_M|F_A^-|^2\,d\mu,\qquad
 J_+=\|F_A^+-(F_A^+)_H\|_2^2.
\end{gathered}
\end{equation}
The Hermitian pairing is conjugate-linear in its first entry.  The
maximum principle gives $0\le\rho\le4$, hence $|\Phi|\le2$.
Subscript $H$ denotes harmonic projection.  This definition of $J_+$
is independent of the harmless sign convention for the imaginary
representative of the Chern class.

At $\beta=2$ the reference weight and spinor amplitude are
\begin{equation}\label{eq:lc-reference}
 c=\sqrt6,\qquad \rho_0=\frac4{\sqrt6},\qquad
 \omega=\frac{\sqrt3}{2}\psi,\qquad
 v=|\omega|^2=\frac34\rho^2,\qquad
 W_0=\frac12\left(\frac v3 I-\omega\otimes\omega\right).
\end{equation}
In particular $c\rho_0=4$.  The expression $W_0$ is a globally smooth
polynomial in $\Phi$, including at its zeros.

\begin{lemma}[Retained estimates for one actual solution]\label{lem:lc-retained}
The majorant in \eqref{eq:lc-sw} may be selected so that the same solution
satisfies
\begin{equation}\label{eq:lc-retained}
\begin{split}
 J_++\int_M\bigl(&|\rho-\rho_0|^2+|f-c|^2+|U|^2+|d\rho|^2
       +|\operatorname{Im}\zeta|^2\\
       &+|W^+-W_0|^2+|\Ric+g|^2\bigr)\,d\mu\le C\Theta.
\end{split}
\end{equation}
It also satisfies
\begin{equation}\label{eq:lc-low-and-slack}
 \int_{\{\rho\le1\}}|\nabla_A\Phi|^2\,d\mu\le C\Theta,
 \qquad |8\sqrt6K-\Sigma|\le C\Theta.
\end{equation}
The constants are universal.  No comparison between solutions at
different weights is involved.
\end{lemma}

\begin{proof}
We first record precisely the estimates from Section~\ref{sec:weighted-pencil}
that are used.  Set $u=3\rho/2$ and $e=s+4$.  Its deficit decomposition,
moment estimate, and retained nonnegative derivative terms give
\begin{equation}\label{eq:lc-moment-input}
\begin{gathered}
 \int_M|q-6|^2\,d\mu\le C\Xi,\qquad
 \int_M e^2\,d\mu\le4rV,\qquad
 \int_M Q(u-\sqrt Q)^2\,d\mu\le C\Theta,\\
 D_0\le C\Theta,\qquad
 0\le C_{\mathrm{act}}-72\pi^2p\le C\Theta,
 \qquad C_{\mathrm{act}}=\int_M f^2u^2\,d\mu,
\end{gathered}
\end{equation}
where the derivative remainder is
\begin{equation}\label{eq:lc-D0}
\begin{split}
 D_0={}&2\int_M|d\rho|^2\,d\mu
 +\int_M(4\rho|\nabla_A\Phi|^2-|\nabla\psi|^2)\,d\mu\\
 &+2\int_M|d\psi|^2\,d\mu+2\int_M E_\lambda\,d\mu,
 \qquad
 E_\lambda=\langle W^+\psi,\psi\rangle+\frac w3\rho^2\ge0.
\end{split}
\end{equation}
These are estimates for the actual solution of \eqref{eq:lc-sw} with
the chosen smooth majorant.  All its approximation terms are included
in $\varepsilon V$.

Clipping at $3$ does not increase the distance to $6$, so
$\|\sqrt Q-c\|_2^2\le C\Xi$.  Since $Q\ge3$, the moment estimate
also bounds $\|u-\sqrt Q\|_2^2$ by $C\Theta$.  Their sum controls
$\|\rho-\rho_0\|_2^2$.  The elementary clipping estimate
\[
 (f_0-\sqrt Q)^2\le f_0^2-Q\le\frac9{16}e^2
\]
then gives $\|f-c\|_2^2\le C\Theta$.  The potential is
$U=e+c(\rho-\rho_0)+(f-c)\rho$, so $\rho\le4$ gives its claimed
$L^2$ bound.  Direct differentiation of the quadratic map gives
\begin{equation}\label{eq:lc-hopf}
 d\rho=2\operatorname{Re}\zeta,\qquad
 |\nabla\psi|^2=4\rho|\nabla_A\Phi|^2-4|\operatorname{Im}\zeta|^2.
\end{equation}
Consequently \eqref{eq:lc-D0} controls both displayed first-order errors.

For the full Weyl error one retains the exact square identity
\begin{equation}\label{eq:lc-weyl-square}
 |W^+-W_0|^2
 =\frac16\left(w-\frac34\rho^2\right)^2+\frac34E_\lambda.
\end{equation}
Indeed $W_0=(3/8)(\rho^2 I/3-\psi\otimes\psi)$.  The deficit
decomposition controls $w-2$ in $L^2$, and the bounded amplitude
estimate controls $3\rho^2/4-2$.  Thus \eqref{eq:lc-weyl-square}
is bounded by $C\Theta$ after integration.  Finally Hodge projection
and the curvature equation give
\begin{equation}\label{eq:lc-two-retained-identities}
 4K+\int_M s\rho\,d\mu+\int_M f\rho^2\,d\mu=0,
 \qquad
 \int_M f^2\rho^2\,d\mu=32\pi^2p+8J_+.
\end{equation}
In particular $8J_+=(4/9)(C_{\mathrm{act}}-72\pi^2p)$, completing
\eqref{eq:lc-retained}.

To prove the low-amplitude estimate, fix $1<b<4/3$ and a smooth
nonincreasing function $\chi$ equal to $1$ on $[0,1]$ and to $0$
on $[b,\infty)$.  The amplitude estimate gives
$\Vol\{\rho<b\}\le C\Theta$.  Testing the last equation in
\eqref{eq:lc-sw} against $\chi(\rho)\Phi$ gives the exact identity
\begin{equation}\label{eq:lc-local-spinor-test}
 4\int_M\chi|\nabla_A\Phi|^2\,d\mu
 +2\int_M\chi'|d\rho|^2\,d\mu
 +\int_M\chi U\rho\,d\mu=0.
\end{equation}
The second integral is bounded in absolute value by $C\Theta$.
The last one is at most
$b\|U\|_2\Vol\{\rho<b\}^{1/2}\le C\Theta$.  This proves the
first estimate in \eqref{eq:lc-low-and-slack} without dividing by
$\rho$ at a zero.

For every constant $c>0$, expansion of a square and substitution of
\eqref{eq:lc-two-retained-identities} give
\begin{equation}\label{eq:lc-exact-slack}
 \Sigma+\int_M(f-c)^2\rho^2\,d\mu
 =\|s+c\rho\|_2^2+8cK+8J_+.
\end{equation}
The coefficient of $\rho^2$ after the scalar--spinor cross term is
eliminated is $f^2-2cf+c^2$.  Thus the term on the left is retained
for a variable weight.  With $c=\sqrt6$, the identity
$s+c\rho=e+c(\rho-\rho_0)$ and \eqref{eq:lc-retained} bound all
terms in \eqref{eq:lc-exact-slack} except $8cK$ and $\Sigma$ by
$C\Theta$.  This proves the second estimate.  No derivative of $f$
has been taken.
\end{proof}

\subsection{The curvature model and the refined Hessian estimate}

\begin{lemma}[Hessian estimate with retained background action]\label{lem:lc-hessian}
The actual solution of Lemma~\ref{lem:lc-retained} satisfies
\begin{equation}\label{eq:lc-refined-hessian}
 \|\nabla_A^2\Phi\|_2^2+\|\nabla_A\Phi\|_4^4
 \le C(K+\Theta+B).
\end{equation}
\end{lemma}

\begin{proof}
Let $\mathcal F$ be the curvature of the connection on the positive
spinor bundle.  Define an algebraic comparison curvature $\mathcal F_b$
by inserting
\begin{equation}\label{eq:lc-background}
 s_b=-4,\qquad \Ric_b^\circ=0,\qquad W_b^+=W_0,
 \qquad F_{A,b}^+=ic\sigma(\Phi),\qquad F_{A,b}^-=0
\end{equation}
into the positive-spinor curvature decomposition.  It is used as an
algebraic tensor; it need not be the curvature of another connection.
The $W^-$ block does not act on positive spinors.  Since $\rho\le4$,
this comparison satisfies $|\mathcal F_b|\le C$, and
\eqref{eq:lc-retained} gives
\begin{equation}\label{eq:lc-background-error}
 \|\mathcal F-\mathcal F_b\|_2^2\le C(\Theta+B).
\end{equation}
For example, the determinant self-dual difference is
$i(f-c)\sigma(\Phi)$ and is controlled by the retained weight
estimate; its anti-self-dual difference contributes exactly $B$.

The action of $\mathcal F_b$ is smaller than its norm.  In a compatible
self-dual frame with its first axis along $\psi$, the eigenvalues of
$W_0$ are $-\rho^2/4,\rho^2/8,\rho^2/8$.  In the curvature and
Clifford conventions of \eqref{eq:lc-sw}, the axial and transverse
coefficients of $\mathcal F_b\Phi$, up to fixed common normalization
factors, are
\begin{equation}\label{eq:lc-annihilation-factors}
\begin{split}
 \frac{\rho^2}{4}+\frac13-\frac{c\rho}{4}
 &=\frac14\left(\rho-\frac4{\sqrt6}\right)
               \left(\rho-\frac2{\sqrt6}\right),\\
 \frac{\rho^2}{8}-\frac13
 &=\frac18\left(\rho-\frac4{\sqrt6}\right)
               \left(\rho+\frac4{\sqrt6}\right).
\end{split}
\end{equation}
Both vanish at $\rho_0$.  Boundedness of $\rho$ therefore proves
\begin{equation}\label{eq:lc-background-annihilates}
 |\mathcal F_b\Phi|^2\le C|\rho-\rho_0|^2,
 \qquad \|\mathcal F\Phi\|_2^2\le C(\Theta+B).
\end{equation}
The estimate extends through $\Phi=0$ because the expressions are
polynomial.  No derivative of $\mathcal F_b$ is used here.

Write $\Pi_i=\nabla_{A,i}\Phi$, $H_{ij}=\nabla_{A,i}\nabla_{A,j}\Phi$
with its Levi--Civita correction, and $L=\operatorname{tr}H$.
The exact covariant identity is
\begin{equation}\label{eq:lc-exact-hessian}
\begin{split}
 \|H\|_2^2={}&\|L\|_2^2
 -\int_M\Ric_{ij}\operatorname{Re}\langle\Pi_i,\Pi_j\rangle\,d\mu\\
 &+\sum_{i,j}\int_M\operatorname{Re}
                  \langle\mathcal F_{ij}\Pi_j,\Pi_i\rangle\,d\mu
 +\frac12\sum_{i,j}\|\mathcal F_{ij}\Phi\|_2^2.
\end{split}
\end{equation}
To see why no curvature derivative appears, first use
$H_{ij}-H_{ji}=\mathcal F_{ij}\Phi$ to write
\[
 \|H\|_2^2
 =\sum_{i,j}\int_M\operatorname{Re}\langle H_{ij},H_{ji}\rangle\,d\mu
  +\frac12\sum_{i,j}\|\mathcal F_{ij}\Phi\|_2^2.
\]
Integrate the mixed Hessian term once, commute derivatives on the
bundle-valued one-form $\Pi$, and integrate the resulting trace term
once more.  The cotangent curvature contracts to Ricci, and the bundle
curvature acts on $\Pi$, giving \eqref{eq:lc-exact-hessian}.  The
sums displayed there are double sums, whereas exterior two-form norms
use $i<j$; the final factor $1/2$ accounts for this distinction.

The bounded-section interpolation used in
Section~\ref{sec:intrinsic-calibration} also follows directly here.
Integration of the divergence of
$|\Pi|^2\operatorname{Re}\langle\Phi,\Pi\rangle^\sharp$, followed
by $|L|\le2|H|$, yields
\[
 \int_M|\Pi|^4\,d\mu
 \le4\|\Phi\|_\infty\int_M|H||\Pi|^2\,d\mu.
\]
Thus
\begin{equation}\label{eq:lc-interpolation}
 \|\Pi\|_4^2\le4\|\Phi\|_\infty\|H\|_2\le8\|H\|_2.
\end{equation}
This uses integration by parts, not a Sobolev inequality.

Equation~\eqref{eq:lc-sw} gives $L=U\Phi/4$ and hence
$\|L\|_2^2\le C\Theta$.  In \eqref{eq:lc-exact-hessian}, split
$\Ric=-g+(\Ric+g)$ and
$\mathcal F=\mathcal F_b+(\mathcal F-\mathcal F_b)$.
The two bounded background contributions cost at most $CK$.
H\"older's inequality and \eqref{eq:lc-interpolation} bound the
remaining terms quadratic in $\Pi$ by
$C\sqrt{\Theta+B}\,\|H\|_2$.  The last term in
\eqref{eq:lc-exact-hessian} is bounded by
\eqref{eq:lc-background-annihilates}.  Therefore
\[
 \|H\|_2^2\le C(K+\Theta+B)
                 +C\sqrt{\Theta+B}\,\|H\|_2.
\]
Young's inequality absorbs the last factor into the left side.
Equation~\eqref{eq:lc-interpolation} then proves
\eqref{eq:lc-refined-hessian}.  In particular no additive multiple
of $V$ remains.
\end{proof}

\subsection{Class-square absorption and the weak Bianchi bootstrap}

The determinant and chiral Chern--Weil identities have the following
useful exact form:
\begin{equation}\label{eq:lc-chern-identities}
 B=4\pi^2(p-a^2)+J_+,
 \qquad
 36\pi^2(p-T)=\frac92P+\frac98\Sigma-D.
\end{equation}
For the second identity, substitute \eqref{eq:lc-energies} into
\[
 4\pi^2T=\int_M\left(2|W^+|^2+\frac{s^2}{24}
                              -\frac12|\Ric^\circ|^2\right)d\mu.
\]
Consequently, for $Y=(T-a^2)_+$,
\begin{equation}\label{eq:lc-absorb-B}
\begin{split}
 B&\le\frac12P+\frac18\Sigma-\frac19D+J_++4\pi^2Y\\
  &\le\sqrt6K+C\Theta+4\pi^2Y.
\end{split}
\end{equation}
We used $P\le rV$, $D\ge0$, and
\eqref{eq:lc-low-and-slack}.  This estimates $B$ by an as yet unknown
spinor energy; it does not assume that $B$ is small.  Lemma~\ref{lem:lc-hessian}
now gives
\begin{equation}\label{eq:lc-hessian-after-square}
 \|\nabla_A^2\Phi\|_2^2+\|\nabla_A\Phi\|_4^4
 \le C(K+\Theta+Y).
\end{equation}

Let $\mathcal W(W_0)$ be the Weyl four-tensor corresponding to the
endomorphism $W_0$, and put
$\mathfrak j=\operatorname{div}\mathcal W(W_0)$.  The weak contracted
Bianchi identity from Section~\ref{sec:intrinsic-calibration}, applied
to any smooth test tensor $Z$, says
\begin{equation}\label{eq:lc-weak-bianchi}
 \left|\int_M\langle\mathfrak j,Z\rangle\,d\mu\right|
 \le C\bigl(\|W^+-W_0\|_2+\|\Ric+g\|_2\bigr)\|\nabla Z\|_2.
\end{equation}
Indeed, split the divergence into the difference from the actual
self-dual Weyl tensor and its actual divergence.  Integrate the former
by parts.  The latter is the self-dual projection of a fixed linear
combination of derivatives of Ricci and scalar curvature.  Since
$\nabla g=0$ and $s+4=\operatorname{tr}(\Ric+g)$, another integration
by parts gives precisely the right side of
\eqref{eq:lc-weak-bianchi}.  No estimate of a derivative of Ricci
is assumed.

The polynomial $W_0$ has degree four in the bounded spinor.  Since its
defining algebraic maps are parallel,
\[
 |\nabla^2W_0|\le C\bigl(|\nabla_A^2\Phi|+|\nabla_A\Phi|^2\bigr),
 \qquad |\nabla\mathfrak j|\le C|\nabla^2W_0|.
\]
No derivative commutation is required.  Taking $Z=\mathfrak j$ in
\eqref{eq:lc-weak-bianchi} and using
\eqref{eq:lc-hessian-after-square} gives
\begin{equation}\label{eq:lc-bianchi-upper}
 \|\mathfrak j\|_2^2\le C\sqrt{\Theta(K+\Theta+Y)}.
\end{equation}

The same test controls $K$ from below.  The polynomial divergence
contraction in Section~\ref{sec:intrinsic-calibration}, specialized
to \eqref{eq:lc-reference}, is
\begin{equation}\label{eq:lc-divergence-exact}
 |\mathfrak j|^2
 =\frac14v|\nabla\omega|^2+\frac12v|d\omega|^2
       -\frac1{12}|dv|^2
       +\frac12\langle\iota_{\nabla v}\omega,*d\omega\rangle.
\end{equation}
It is valid for an arbitrary smooth self-dual form, including at a
zero.  The last term is bounded in absolute value by
$\tfrac12\sqrt v\,|dv||d\omega|$.  Young's inequality therefore gives
\begin{equation}\label{eq:lc-divergence-coercive}
 \frac14v\bigl(|\nabla\omega|^2+|d\omega|^2\bigr)
 \le|\mathfrak j|^2+\frac13|dv|^2.
\end{equation}
On $\{\rho>1\}$, equations~\eqref{eq:lc-reference} and
\eqref{eq:lc-hopf} read
\[
 |\nabla\omega|^2=3\rho|\nabla_A\Phi|^2
                         -3|\operatorname{Im}\zeta|^2,
 \qquad dv=\frac32\rho\,d\rho.
\]
Here $v\ge3/4$.  Integrating \eqref{eq:lc-divergence-coercive} and
using the retained current and amplitude derivatives gives
\[
 \int_{\{\rho>1\}}|\nabla_A\Phi|^2\,d\mu
 \le C\bigl(\|\mathfrak j\|_2^2+\Theta\bigr).
\]
Adding the low-amplitude estimate in
\eqref{eq:lc-low-and-slack} proves
\begin{equation}\label{eq:lc-bianchi-lower}
 K\le C\bigl(\Theta+\|\mathfrak j\|_2^2\bigr).
\end{equation}
The low-amplitude term is essential: no inverse power of $\rho$ has
been extended across its zero set.

Combining \eqref{eq:lc-bianchi-upper} and
\eqref{eq:lc-bianchi-lower}, and separating the three terms under
the square root, gives
\[
 K\le C\Theta+C\sqrt{\Theta K}+C\sqrt{\Theta Y}.
\]
Absorbing the middle term by Young's inequality yields
\begin{equation}\label{eq:lc-K-closed}
 K\le C\bigl(\Theta+\sqrt{\Theta Y}\bigr).
\end{equation}

\begin{theorem}[Normalized linear coercivity]\label{thm:lc-normalized}
Suppose the actual input above satisfies $\kappa\le1$.  For every
sufficiently small selected allowance $\varepsilon>0$, the same
weighted solution satisfies
\begin{equation}\label{eq:lc-shortfall-fields}
\begin{split}
 K+\Sigma+(p-T)_+&\le C\bigl(\Theta+\sqrt{\Theta Y}\bigr),\\
 B+\|\nabla_A^2\Phi\|_2^2+\|\nabla_A\Phi\|_4^4
       +\int_M|\zeta|^2|\nabla_A\Phi|^2\,d\mu
       &\le C(\Theta+Y).
\end{split}
\end{equation}
In particular, if $a^2\ge T$, every term displayed is bounded by
$C\Theta$, and the solution-independent quantities obey
\begin{equation}\label{eq:lc-normalized-topology}
 0\le\Sigma\le C(rV+D),\qquad
 0\le p-T\le C_T(rV+D).
\end{equation}
Without that square condition they obey
\begin{equation}\label{eq:lc-normalized-shortfall}
 \Sigma+(p-T)_+
 \le C\bigl(rV+D+\sqrt{(rV+D)Y}\bigr).
\end{equation}
\end{theorem}

\begin{proof}
Equation~\eqref{eq:lc-K-closed} and the slack comparison in
\eqref{eq:lc-low-and-slack} bound $K$ and $\Sigma$ as stated.
The second identity in \eqref{eq:lc-chern-identities}, with
$P\le rV\le\Theta$ and $D\ge0$, gives the bound for $(p-T)_+$.
Equations~\eqref{eq:lc-absorb-B} and
\eqref{eq:lc-hessian-after-square}, together with
$\sqrt{\Theta Y}\le(\Theta+Y)/2$, give the second line.
Finally $|\zeta|^2\le\rho|\nabla_A\Phi|^2\le4|\nabla_A\Phi|^2$
controls its mixed term.  If $a^2\ge T$, then
$p=a^2+\|(a_g)^-\|^2\ge T$, so the second quantity in
\eqref{eq:lc-normalized-topology} is nonnegative.  Sending the freely
selected allowance to zero for the fixed metric and fixed class proves
the solution-independent assertions.  No convergence of the selected
solutions is required.
\end{proof}

\begin{remark}\label{rem:lc-square-cost}
The shortfall has two different effects.  The scalar and topological
conclusions pay $\sqrt{\Xi Y}$, whereas the determinant, Hessian,
and mixed fourth-order energies can pay $Y$ itself.  The first line of
\eqref{eq:lc-shortfall-fields} cannot be substituted for the second
when estimating the mixed current.  In the square-controlled case,
the Chern identity $4c_2(S^+)[M]=a^2-T$ also gives
$c_2(S^+)[M]\le C\Xi$; its positivity is not needed for the analytic
bootstrap.
\end{remark}

There is a useful nearby-parameter consequence which preserves the
same actual solution.  For $2\le\beta\le3$, let
$\delta=\beta-2$ and $\Xi_\beta=rV+D_\beta$.  The exact decomposition
from Section~\ref{sec:weighted-pencil} is
\[
 D_\beta=A_0+\frac1\beta\int_M(s+\beta w)^2\,d\mu,
 \qquad A_0\ge0.
\]
Writing $z=s+\beta w$ gives
\[
 D_2=A_0+\frac12\int_M(z-\delta w)^2\,d\mu
       \le\beta D_\beta+\delta^2\int_Mw^2\,d\mu.
\]
Since $w=(z-(s+4)+4)/\beta$, successive squared triangle
inequalities imply
\[
 \int_Mw^2\,d\mu
 \le\frac4\beta D_\beta+\frac{16}{\beta^2}rV
                         +\frac{32}{\beta^2}V
 \le2D_\beta+4rV+8V.
\]
Consequently
\begin{equation}\label{eq:lc-nearby-pencil}
 \Xi_2\le5\Xi_\beta+8(\beta-2)^2V.
\end{equation}
Whenever $\Xi_2/V\le1$, Theorem~\ref{thm:lc-normalized} applies
with this bound.  The majorant \eqref{eq:lc-weight} is independent
of $\beta$, so these inequalities compare deficits of one chosen
weighted solution.  This observation concerns $\beta$ near $2$.

\subsection{A global relative-defect window}

We next show how the normalized theorem gives an all-metric statement,
including metrics outside its small-error range.  Suppose that for
every smooth metric on $M$ one can choose actual input as above with
\begin{equation}\label{eq:lc-all-metric-input}
 p\ge A>0,\qquad T_+=A(1-h),\qquad 0<h\le\frac13.
\end{equation}
The constant $A$ and the topological parameter $h$ are fixed for the
manifold.  Write
\begin{equation}\label{eq:lc-global-notation}
 U_{\mathrm{top}}=\frac{3\min\{T_+,T_-\}}A,\qquad
 e=\frac{R}{8\pi^2A},\qquad
 \E(g)=\inf_{\lambda\ne0}
       \frac{\int_M|\Ric-\lambda g|^2\,d\mu}{\lambda^2V}.
\end{equation}
With $I=\int_Ms\,d\mu$ and
$L_s=\int_M(s-I/V)^2\,d\mu$, minimizing the quadratic polynomial in
$1/\lambda$ gives
\begin{equation}\label{eq:lc-E-formula}
 \E=4-\frac{I^2}{VR}=\frac{4P+L_s}{R}.
\end{equation}
The formula remains valid when $I=0$, in which case the infimum is
$4$ and need not be attained at a finite target.  The two chiral
Chern--Weil identities give
\begin{equation}\label{eq:lc-chiral-E}
 e(1-\E)\le U_{\mathrm{top}}.
\end{equation}
For example,
$R-4P=24\pi^2T_+-12\int_M|W^+|^2\,d\mu$; the analogous formula
for $W^-$ and $\E\ge4P/R$ prove \eqref{eq:lc-chiral-E}.

\begin{theorem}[Linear relative-defect window]\label{thm:lc-E-window}
Under \eqref{eq:lc-all-metric-input}, let
$y=(T_+-a^2)_+/A$ for the class selected at $g$.  A universal $c_*>0$
satisfies
\begin{equation}\label{eq:lc-E-shortfall-window}
 \E(g)+(U_{\mathrm{top}}-1-h)_+
       \ge c_*\frac{h^2}{h+y}.
\end{equation}
In particular, $a^2\ge T_+$ gives the lower bound $c_*h$.
More generally, $y\le Lh$ gives $c_*h/(1+L)$.
\end{theorem}

\begin{proof}
Set $\Delta=(U_{\mathrm{top}}-1-h)_+$.  If
$U_{\mathrm{top}}\le0$, equation~\eqref{eq:lc-chiral-E} implies
$\E\ge1$, which suffices because $h^2/(h+y)\le h\le1/3$.
Otherwise $0<U_{\mathrm{top}}\le3$.  If $\E+\Delta\ge1/49$,
then $\E+\Delta\ge3h/49$, again sufficient after fixing a smaller
universal constant.  It remains to treat
\begin{equation}\label{eq:lc-small-global-branch}
 \E+\Delta<\frac1{49}.
\end{equation}

The sign of the scalar average in this branch follows from the actual
input.  If $I\ge0$, then on $\{s<0\}$ one has
$s_-^2\le(s-I/V)^2$.  The negative scalar estimate and
\eqref{eq:lc-E-formula} imply
\[
 32\pi^2A\le\int_Ms_-^2\,d\mu\le L_s\le R\E,
 \qquad 4\le e\E.
\]
Since $\E<1$, combining this with \eqref{eq:lc-chiral-E} yields
$4(1-\E)\le U_{\mathrm{top}}\E$, hence
$\E\ge4/(U_{\mathrm{top}}+4)\ge4/7$, contradicting
\eqref{eq:lc-small-global-branch}.  Thus $I<0$.

Rescale the metric by the positive constant $-R/I$.  The minimizing
negative target becomes $-g$, and the rescaled quantities satisfy
\begin{equation}\label{eq:lc-optimal-rescaling}
 r=\E,\qquad V=\frac{R}{4-\E}.
\end{equation}
The integrals $R,P,D$, the projection $p$, and the class square are
unchanged by this constant rescaling in dimension four.  Relabel the
rescaled metric as $g$, and put $z=p/A-1\ge0$.  Direct elimination
of the Weyl energy gives
\begin{equation}\label{eq:lc-D-energy}
 \frac DA=36\pi^2\{e-(1+h)-2z\}.
\end{equation}
It follows that $e\ge1+h$, while
$e\le U_{\mathrm{top}}/(1-\E)\le6$.  More precisely,
\[
 e-(1+h)\le\frac{\Delta+(1+h)\E}{1-\E}
               \le2\Delta+\frac83\E.
\]
Therefore
\begin{equation}\label{eq:lc-global-error-comparison}
 \frac DA\le72\pi^2\Delta+96\pi^2\E,
 \qquad
 \frac{rV}{A}=\frac{8\pi^2e\E}{4-\E}
                     \le\frac{96}{7}\pi^2\E.
\end{equation}
In particular
\begin{equation}\label{eq:lc-Xi-global}
 \frac\Xi A\le112\pi^2(\E+\Delta),\qquad
 \kappa\le49(\E+\Delta)<1.
\end{equation}
For the second inequality use
$A/V=(4-\E)/(8\pi^2e)\le1/(2\pi^2)$ directly in the first
inequality of \eqref{eq:lc-global-error-comparison}; it gives
$D/V\le36\Delta+48\E$, and $r=\E$ supplies the last term.

Now $p-T_+=A(h+z)\ge Ah$.  The normalized shortfall estimate
\eqref{eq:lc-normalized-shortfall}, after the allowance has been
sent to zero, gives, with $x=\Xi/A$,
\[
 h\le C\bigl(x+\sqrt{xy}\bigr).
\]
If $x\ge h$, then $x\ge h^2/(h+y)$.  If $x<h$, squaring yields
$h^2\le2C^2x(x+y)\le2C^2x(h+y)$.  Thus
$x\ge c h^2/(h+y)$ in either case.  Combining with
\eqref{eq:lc-Xi-global} proves the theorem on the small branch.
The two preceding branches complete its all-metric quantifier.
\end{proof}

\subsection{Conformal transfer without scalar variance}

\begin{theorem}[Scale-invariant scalar and topological coercivity]\label{thm:lc-scalar-free}
For every metric with the actual conformal supply specified above,
without a small-error assumption, there are universal constants
$C,C_T>0$ such that
\begin{equation}\label{eq:lc-scalar-free}
 \Sigma\le C\bigl(P+D+\sqrt{(P+D)Y}\bigr),\qquad
 (p-T)_+\le C_T\bigl(P+D+\sqrt{(P+D)Y}\bigr).
\end{equation}
If $a^2\ge T$, this reduces to
\begin{equation}\label{eq:lc-scalar-free-square}
 \Sigma\le C(P+D),\qquad p-T\le C_T(P+D).
\end{equation}
\end{theorem}

\begin{proof}
The conformal class has negative Yamabe constant.  Otherwise its
constant-scalar-curvature representative would have $s\ge0$.  For
the actual ordinary monopole supplied on that representative, the
integrated spinor equation would force the spinor and its self-dual
curvature to vanish.  This contradicts $p>0$.  Let $\widehat g$
be the negative Yamabe representative rescaled to
$\widehat s=-4$.  The same Spin$^c$ structure is available there.
The Hodge star on two-forms is conformally invariant in dimension
four, so the harmonic decomposition of $a$, the value of $p$, and
$Y$ are unchanged.

Write $g=u^2\widehat g$, $u>0$, and use the nonnegative Laplacian
$\widehat\Delta=d^*d$.  The scalar transformation law is
$s_g u^3=6\widehat\Delta u-4u$.  Thus
\[
 S_2=16\widehat V-48\int_Mu^{-1}\widehat\Delta u\,d\widehat\mu
             +36\int_M(u^{-1}\widehat\Delta u)^2\,d\widehat\mu.
\]
Since
$\int_Mu^{-1}\widehat\Delta u\,d\widehat\mu
 =-\int_Mu^{-2}|du|_{\widehat g}^2\,d\widehat\mu$, it follows that
\begin{equation}\label{eq:lc-conformal-delta}
 \delta:=S_2-\widehat S_2
 =48\int_Mu^{-2}|du|_{\widehat g}^2\,d\widehat\mu
  +36\int_M(u^{-1}\widehat\Delta u)^2\,d\widehat\mu\ge0.
\end{equation}
The integrated Weyl energy is conformally invariant.  Chiral
Chern--Weil and \eqref{eq:lc-energies} therefore give the exact
comparisons
\begin{equation}\label{eq:lc-conformal-energies}
 P-\widehat P=\frac\delta{12},\qquad
 D-\widehat D=\frac32\delta,\qquad
 \Sigma-\widehat\Sigma=\delta.
\end{equation}
These are energy identities and do not require a small conformal factor.

Put $\Pi=P+D$ and $\widehat\Pi=\widehat P+\widehat D$.
On $\widehat g$, the negative target residual satisfies
$\widehat r\widehat V=\widehat P$.  If
$\widehat\Pi/\widehat V\le1$, Theorem~\ref{thm:lc-normalized}
and its allowance limit give
\[
 \widehat\Sigma\le C_0\bigl(\widehat\Pi+
                                      \sqrt{\widehat\Pi Y}\bigr).
\]
If $\widehat\Pi/\widehat V>1$, then instead
\[
 0\le\widehat\Sigma=16\widehat V-32\pi^2p
                  \le16\widehat V<16\widehat\Pi.
\]
The same bound holds in both branches after enlarging the universal
constant.  Equation~\eqref{eq:lc-conformal-energies}
gives
$\Pi=\widehat\Pi+(19/12)\delta$.  Hence
$\widehat\Pi\le\Pi$ and $\delta\le12\Pi/19$, and the first
inequality in \eqref{eq:lc-scalar-free} follows from
$\Sigma=\widehat\Sigma+\delta$.

Finally, the exact topological identity in
\eqref{eq:lc-chern-identities} gives
\[
 36\pi^2(p-T)\le\frac92P+\frac98\Sigma.
\]
The right side is nonnegative; bounding it by the first inequality
proves the assertion for $(p-T)_+$.  If $a^2\ge T$, then $Y=0$
and $p\ge a^2\ge T$, proving \eqref{eq:lc-scalar-free-square}.
\end{proof}

The conclusion is about the original metric's curvature energies.
The proof neither preserves a pointwise Ricci bound under the
conformal change nor identifies the hatted spinor with a weighted
spinor on the original metric.

\subsection{The lower energy wall and the traceless-Ricci window}

For a metric with $p\ge A>0$ and $T=A(1-h)$, where $h>0$, define
\begin{equation}\label{eq:lc-J-definition}
 \J(g)=\frac{4P}{R},\qquad z=\frac pA-1\ge0,
 \qquad y=\frac YA.
\end{equation}
Then $0\le\J\le4$ and $\J\le\E$.  Equations
\eqref{eq:lc-D-energy} and \eqref{eq:lc-J-definition} give
\begin{equation}\label{eq:lc-Pi-reduction}
 \frac PA=2\pi^2e\J,\qquad
 \frac{P+D}{A}
   =36\pi^2\left[e\left(1+\frac\J{18}\right)-1-h-2z\right].
\end{equation}
In particular $D\ge0$ gives $e\ge1+h+2z$.

\begin{theorem}[Lower Ricci-energy wall]\label{thm:lc-energy-wall}
If $a^2\ge T_+$, there is a universal $c_0>0$ such that
\begin{equation}\label{eq:lc-wall-square}
 e\left(1+\frac\J{18}\right)
          \ge1+(1+c_0)h+(2+c_0)z.
\end{equation}
For an arbitrary square shortfall there is a universal $c_s>0$ such that
\begin{equation}\label{eq:lc-wall-shortfall}
 e\left(1+\frac\J{18}\right)
    \ge1+h+2z+c_s\frac{(h+z)^2}{h+z+y}.
\end{equation}
These assertions require $h>0$, but no upper bound on $h$ or $e$ and
no upper energy budget from the opposite chirality.
\end{theorem}

\begin{proof}
In the square-controlled case,
$A(h+z)=p-T\le C_T(P+D)$ by
Theorem~\ref{thm:lc-scalar-free}.  Substitute
\eqref{eq:lc-Pi-reduction} and take
$c_0=(36\pi^2C_T)^{-1}$.  For the general case set
$x=(P+D)/A$ and $v_0=h+z>0$.  The same theorem gives
$v_0\le C_T(x+\sqrt{xy})$.  If $x\ge v_0$ then
$x\ge v_0^2/(v_0+y)$.  Otherwise, squaring shows
$v_0^2\le2C_T^2x(x+y)\le2C_T^2x(v_0+y)$.
Thus $x\ge c v_0^2/(v_0+y)$ in both cases.  Substitution into
\eqref{eq:lc-Pi-reduction} proves the claim after dividing the
constant by $36\pi^2$.
\end{proof}

In the range $0<h\le1/3$, equation~\eqref{eq:lc-wall-square} also
implies $\J+(e-1-h)\ge ch$.  Indeed, if $e-1-h\ge h/2$ this is
immediate.  Otherwise $e<1+3h/2\le3/2$, and
$e-1-h+\J/12\ge c_0h$ follows by dropping $z$.
For an Einstein metric allowed by the input one has $\J=0$, so the
same lower wall gives $e\ge1+(1+c_0)h$.  Ricci-flat metrics have
already been excluded by \eqref{eq:lc-nonnegative}.

The two chiral identities now retain the stronger upper budget
\begin{equation}\label{eq:lc-chiral-J}
 e(1-\J)\le U_{\mathrm{top}}.
\end{equation}
No scalar variance has been discarded in this inequality.

\begin{theorem}[Traceless-Ricci linear window]\label{thm:linear-window}
Under \eqref{eq:lc-all-metric-input} and $a^2\ge T_+$, put
$L_h=1+(1+c_0)h$.  If $U_{\mathrm{top}}>0$, then
\begin{equation}\label{eq:lc-J-explicit}
 \J\ge\frac{(L_h-U_{\mathrm{top}})_+}
                    {L_h+U_{\mathrm{top}}/18}.
\end{equation}
If $U_{\mathrm{top}}\le0$, then $\J\ge1$.  In particular,
\begin{equation}\label{eq:lc-J-linear}
 \J+(U_{\mathrm{top}}-1-h)_+\ge ch.
\end{equation}
Without the square condition the corresponding bound is
\begin{equation}\label{eq:lc-J-shortfall}
 \J+(U_{\mathrm{top}}-1-h)_+
       \ge c\frac{h^2}{h+y}.
\end{equation}
Every assertion remains valid with $\E$ in place of $\J$.
\end{theorem}

\begin{proof}
Suppose first that $U_{\mathrm{top}}>0$ and $\J<1$.
Equations~\eqref{eq:lc-wall-square} and \eqref{eq:lc-chiral-J}
give
\[
 L_h(1-\J)\le U_{\mathrm{top}}(1+\J/18).
\]
Rearrangement proves \eqref{eq:lc-J-explicit}.  For $\J\ge1$
its right side is at most $1$, so the assertion still holds.
If $U_{\mathrm{top}}\le0$, equation~\eqref{eq:lc-chiral-J}
and $e>0$ directly give $\J\ge1$.

For the linear bound put $\Delta=(U_{\mathrm{top}}-1-h)_+$.
If $\Delta\ge c_0h/2$, the result follows from that term alone.
Otherwise $L_h-U_{\mathrm{top}}\ge c_0h/2$.  Since
$0<U_{\mathrm{top}}\le3$ and $h\le1/3$, the denominator in
\eqref{eq:lc-J-explicit} is at most $3/2+c_0/3$.
This proves \eqref{eq:lc-J-linear} in the remaining branch.

For a shortfall, the function $v_0\mapsto v_0^2/(v_0+y)$ is
nondecreasing for $v_0>0$.  Dropping $z$ in
\eqref{eq:lc-wall-shortfall} therefore gives the lower wall
$1+h+c_s h^2/(h+y)$.  Repeat the preceding argument with
$c_0h$ replaced by $c_s h^2/(h+y)$; this quantity is at most
$c_sh$, so the denominator remains uniformly bounded.
This proves \eqref{eq:lc-J-shortfall}, including the branches
$\J\ge1$ and $U_{\mathrm{top}}\le0$.  Finally $\E\ge\J$
proves all relative-defect versions.
\end{proof}

\subsection{The meaning of linear order}

The sharpness consequence separates the analytic lower bound from its
construction input.  Suppose a sequence of actual supplied manifolds has
$h_j\to0$, $a^2\ge T_+$ for its selected classes,
$U_{\mathrm{top},j}=1+h_j+O(h_j^2)$, and independently constructed
metrics with
$\inf_g\E(g)\le4h_j/(1+h_j)$.  Theorem~\ref{thm:linear-window}
then gives, for all sufficiently large $j$,
\begin{equation}\label{eq:lc-sharp-order}
 c_1h_j\le\inf_g\J(g)=\inf_g\E(g)
                     \le\frac{4h_j}{1+h_j}.
\end{equation}
Indeed the $O(h_j^2)$ chiral allowance is absorbed in half of the
linear lower bound, and the upper bound is eventually less than $1$,
so the conformal-infimum reduction in
Section~\ref{sec:variational-geometry} identifies the two infima.
The construction used for the upper bound is logically separate
from the coercivity proof.  Equation~\eqref{eq:lc-sharp-order}
determines the order and does not determine the optimal leading
coefficient over all metrics.  The estimate also concerns the
K\"ahler calibration $\beta=2$; no nondegenerate estimate at the
zero-Weyl end of the pencil is asserted here.
